\chapter*{Список использованных источников}
\addcontentsline{toc}{chapter}{Список использованных источников}
\markboth{Список использованных источников}{}
\begin{thebibliography}{99}

\sloppy

\bibitem{golang}
Официальный сайт языка программирования Go [Электронный ресурс]. — URL: \url{https://go.dev/} (дата обращения: 21.12.2025).

\bibitem{fyne}
Официальный сайт библиотеки Fyne [Электронный ресурс]. — URL: \url{https://fyne.io/} (дата обращения: 21.12.2025).

\bibitem{donovan}
Донован, А. А., Керниган, Б. У. Язык программирования Go. — М.: Вильямс, 2016. — 496 с.

\bibitem{gopsutil}
Shirou, G. gopsutil: библиотека для получения информации о системных процессах и использовании системы [Электронный ресурс] // GitHub. — URL: \url{https://github.com/shirou/gopsutil} (дата обращения: 21.12.2025).

\bibitem{tanenbaum}
Таненбаум, Э., Бос, Х. Современные операционные системы. 4-е изд. — СПб.: Питер, 2015. — 1120 с.

\bibitem{patterson_hennessy}
Паттерсон, Д., Хеннесси, Дж. Архитектура компьютера и проектирование компьютерных систем. 4-е изд. — СПб.: Питер, 2012. — 784 с.

\bibitem{go_concurrency}
Chan, A. Concurrency in Go: Tools and Techniques for Developers. — O'Reilly Media, 2017. — 256 p.

\bibitem{sutter_free_lunch}
Sutter, H. The Free Lunch Is Over: A Fundamental Turn Toward Concurrency in Software [Электронный ресурс] // Dr. Dobb's Journal. — 2005. — URL: \url{http://www.gotw.ca/publications/concurrency-ddj.htm} (дата обращения: 21.12.2025).

\bibitem{golang_crypto}
Go Package Documentation: crypto, crypto/cipher, crypto/aes, crypto/rsa, crypto/ed25519 [Электронный ресурс]. — URL: \url{https://pkg.go.dev/crypto} (дата обращения: 31.03.2026).

\bibitem{xcrypto}
Go x/crypto: дополнительные криптографические алгоритмы (в т.ч. ChaCha20-Poly1305) [Электронный ресурс] // GitHub. — URL: \url{https://github.com/golang/crypto} (дата обращения: 31.03.2026).

\bibitem{nistsha}
NIST FIPS PUB 180-4: Secure Hash Standard (SHS) [Электронный ресурс]. — URL: \url{https://csrc.nist.gov/publications/detail/fips/180/4/final} (дата обращения: 31.03.2026).

\bibitem{nistaes}
NIST FIPS PUB 197: Advanced Encryption Standard (AES) [Электронный ресурс]. — URL: \url{https://csrc.nist.gov/publications/detail/fips/197/final} (дата обращения: 31.03.2026).

\bibitem{nistrsa}
NIST SP 800-56B Rev. 2: Recommendation for Pair-Wise Key Establishment Schemes Using Integer Factorization Cryptography [Электронный ресурс]. — URL: \url{https://csrc.nist.gov/publications/detail/sp/800-56b/rev-2/final} (дата обращения: 31.03.2026).

\bibitem{tukey}
Tukey, J. W. Exploratory Data Analysis. — Addison-Wesley, 1977. — 688 p.

\bibitem{hampel}
Hampel, F. R., Ronchetti, E. M., Rousseeuw, P. J., Stahel, W. A. Robust Statistics: The Approach Based on Influence Functions. — Wiley, 1986. — 536 p.

\bibitem{ppg}
Hoaglin, D. C., Mosteller, F., Tukey, J. W. Understanding Robust and Exploratory Data Analysis. — Wiley, 1983. — 468 p.

\bibitem{drepper_memory}
Drepper, U. What Every Programmer Should Know About Memory [Электронный ресурс] // LWN.net. — 2007. — URL: \url{https://lwn.net/Articles/250967/} (дата обращения: 21.12.2025).

\bibitem{intel_optimization}
Intel 64 and IA-32 Architectures Optimization Reference Manual [Электронный ресурс] // Intel Corporation. — URL: \url{https://www.intel.com/content/www/us/en/developer/articles/technical/intel-sdm.html} (дата обращения: 21.12.2025).

\bibitem{geekbench}
Geekbench - Cross-Platform Benchmark [Электронный ресурс]. — URL: \url{https://www.geekbench.com/} (дата обращения: 21.12.2025).

\bibitem{pprof}
Go pprof - Profiling Go programs [Электронный ресурс] // GitHub. — URL: \url{https://github.com/google/pprof} (дата обращения: 21.12.2025).

\bibitem{go_benchmarking}
Go Blog: Benchmarking in Go [Электронный ресурс]. — URL: \url{https://go.dev/blog/} (дата обращения: 31.03.2026).

\bibitem{gopsutil_docs}
gopsutil documentation: CPU, Memory, Disk, Net, Process [Электронный ресурс]. — URL: \url{https://pkg.go.dev/github.com/shirou/gopsutil/v3} (дата обращения: 31.03.2026).

\end{thebibliography}
